\documentclass[12pt,letterpaper]{article}
\usepackage[utf8]{inputenc}
\usepackage[spanish]{babel}
\usepackage{times}
\usepackage[margin=2.54cm]{geometry}
\usepackage{graphicx}
\usepackage{hyperref}
\usepackage{fancyhdr}
\usepackage{verbatim}
\usepackage{longtable}
\usepackage{booktabs}
\usepackage{xcolor}
\usepackage{listings}

% Code listing style
\lstset{
    basicstyle=\ttfamily\small,
    breaklines=true,
    frame=single,
    backgroundcolor=\color{gray!10}
}

% Header/Footer
\pagestyle{fancy}
\fancyhf{}
\rhead{Manual Técnico - Prototipo Web 3D}
\lhead{Ingeniería Multimedia - UNAD}
\rfoot{Página \thepage}

\title{\textbf{Manual Técnico}\\Prototipo Web 3D Interactivo para Visualización Técnica de Drones}
\author{Alexander Woodcock Salomón}
\date{Diciembre 2025}

\begin{document}

\maketitle
\tableofcontents
\newpage

\section{Introducción}
Este documento técnico detalla la arquitectura, componentes y procedimientos de mantenimiento del prototipo Web 3D para visualización técnica de drones. El sistema implementa más de 65 scripts con arquitectura modular, patrones de diseño profesionales y optimizaciones específicas para WebGL.

\subsection{Alcance del Documento}
\begin{itemize}
    \item Arquitectura del sistema y patrones de diseño
    \item Descripción de todos los sistemas implementados
    \item Procedimientos de compilación y despliegue
    \item Guías de extensión y mantenimiento
\end{itemize}

\section{Arquitectura del Sistema}

\subsection{Stack Tecnológico}
\begin{itemize}
    \item \textbf{Motor:} Unity 6.3 LTS
    \item \textbf{Lenguaje:} C\# 11 (Scripting) / HLSL (Shaders)
    \item \textbf{Plataforma:} WebGL 2.0
    \item \textbf{Compilador:} IL2CPP $\rightarrow$ WebAssembly (WASM)
    \item \textbf{Pipeline Gráfico:} Universal Render Pipeline (URP)
    \item \textbf{UI Framework:} UI Toolkit (USS/UXML)
\end{itemize}

\subsection{Patrones de Diseño Implementados}
\begin{itemize}
    \item \textbf{Singleton/PersistentSingleton:} Para managers globales
    \item \textbf{Event Bus:} Comunicación desacoplada entre sistemas
    \item \textbf{State Machine:} Control de estados de la aplicación
    \item \textbf{ScriptableObjects:} Configuración de datos (DronePartData, AssemblyGuideData)
    \item \textbf{Object Pooling:} Reutilización de objetos para rendimiento
\end{itemize}

\subsection{Diagrama de Arquitectura}
\begin{verbatim}
+------------------+     +------------------+     +------------------+
|   AppStateMachine|<--->|     EventBus     |<--->| SelectionManager |
+------------------+     +------------------+     +------------------+
         |                       |                        |
         v                       v                        v
+------------------+     +------------------+     +------------------+
|   GameManager    |     |   AudioManager   |     |  UIManager       |
+------------------+     +------------------+     +------------------+
         |                       |                        |
         v                       v                        v
+------------------+     +------------------+     +------------------+
| OrbitCamera      |     |  SaveSystem      |     | PartCatalogUI    |
+------------------+     +------------------+     +------------------+
\end{verbatim}

\section{Estructura del Proyecto}

\begin{verbatim}
Assets/Scripts/
+-- Core/
|   +-- Data/           # ScriptableObjects (DronePartData)
|   +-- Events/         # Sistema de eventos (EventBus)
|   +-- Managers/       # 25+ Managers (Singletons)
|   +-- Content/        # Componentes de piezas
|   +-- Utils/          # Utilidades (Singleton, TweenEngine)
+-- UI/
|   +-- UIManager.cs    # Control principal UI
|   +-- PartCatalogUI   # Catálogo de piezas
|   +-- EnhancedInfoPanel # Panel de información
|   +-- AssemblyStepUI  # Guía de ensamblaje
+-- Graphics/
|   +-- Shaders/        # Shaders personalizados
|   +-- PostProcessing/ # Efectos de post-procesado
\end{verbatim}

\section{Sistemas Principales}

\subsection{Core Managers (25+ Scripts)}

\subsubsection{AppStateMachine.cs}
Máquina de estados principal. Estados: Loading, Exploration, ExplodedView, PartFocus, Settings, Menu.
\begin{lstlisting}[language=C]
AppStateMachine.Instance.TransitionTo(AppState.ExplodedView);
\end{lstlisting}

\subsubsection{EventBus.cs}
Sistema de eventos pub/sub para comunicación desacoplada.
\begin{lstlisting}[language=C]
EventBus.Publish(new PartSelectedEvent(partData));
EventBus.Subscribe<PartSelectedEvent>(OnPartSelected);
\end{lstlisting}

\subsubsection{SelectionManager.cs}
Gestiona selección de piezas con raycast, integra HighlightSystem, CursorManager y AnalyticsManager.

\subsubsection{ViewModeManager.cs}
Controla 7 modos de visualización:
\begin{itemize}
    \item Realistic - Material original
    \item XRay - Semi-transparente con fresnel
    \item Blueprint - Estilo técnico
    \item SolidColor - Color sólido configurable
    \item Wireframe - Solo líneas
    \item Ghosted - Fantasma
    \item Thermal - Vista térmica
\end{itemize}

\subsubsection{PartCatalogManager.cs}
Búsqueda y filtrado de piezas por nombre, tipo, categoría y material.

\subsubsection{PartVisibilityManager.cs}
Control de visibilidad: ocultar, mostrar, aislar piezas con animaciones de fade.

\subsubsection{CrossSectionManager.cs}
Cortes transversales en ejes X, Y, Z con plano visual y posición configurable.

\subsubsection{DroneStateController.cs}
Control de estado del dron: Off, StartingUp, Idle, Flying, ShuttingDown. Incluye animación de propelas, luces de estado y partículas.

\subsubsection{ModularPartsSystem.cs}
Sistema de slots para intercambio de piezas modulares con animaciones.

\subsection{Herramientas de Ensamblaje}

\subsubsection{AssemblyGuideManager.cs}
Guía paso a paso con:
\begin{itemize}
    \item Pasos secuenciales con navegación
    \item Herramientas requeridas por paso
    \item Tiempos estimados
    \item Niveles de dificultad (1-5)
    \item Advertencias de seguridad
    \item Tips de instalación
\end{itemize}

\subsubsection{MeasurementTool.cs}
Herramienta de medición con modos: Distancia, Ángulo, Radio.

\subsubsection{ConnectionPointsViewer.cs}
Visualización de puntos de conexión (tornillos, snaps, soldaduras, cables).

\subsubsection{BillOfMaterialsManager.cs}
Generación de lista de materiales con exportación CSV.

\subsubsection{AnnotationSystem.cs}
Sistema de notas 3D con billboard text.

\subsubsection{AssemblyChecklist.cs}
Lista de verificación de piezas con progreso.

\subsection{Sistemas de Audio y Visual}

\subsubsection{AudioManager.cs}
Control de audio con volúmenes separados (Master, SFX, Music) y persistencia.

\subsubsection{QualityManager.cs}
Ajuste dinámico de calidad basado en FPS.

\subsubsection{TweenEngine.cs}
Motor de animaciones ligero (alternativa a DOTween) con 8 tipos de easing.

\subsubsection{NotificationManager.cs}
Sistema de notificaciones toast animadas.

\section{ScriptableObjects}

\subsection{DronePartData}
Configuración completa de cada pieza:
\begin{itemize}
    \item Info General: nombre, ID, tipo, descripción
    \item Specs Técnicos: peso, dimensiones, función
    \item Materiales: tipo, propiedades, fabricante
    \item Vista Explosionada: dirección, distancia, prioridad
    \item Info de Ensamblaje: herramientas, torque, tornillos, dificultad
\end{itemize}

\subsection{AssemblyGuideData}
Definición de guías de ensamblaje con lista de pasos configurables.

\section{Compilación y Despliegue}

\subsection{Requisitos de Build}
\begin{enumerate}
    \item Unity 6.3 LTS
    \item WebGL Build Support instalado
    \item URP configurado
\end{enumerate}

\subsection{Configuración de Build}
\begin{enumerate}
    \item File $>$ Build Settings $>$ WebGL
    \item Player Settings:
    \begin{itemize}
        \item Compression: Brotli (producción) o Disabled (desarrollo)
        \item Code Stripping: High
        \item IL2CPP Code Generation: Faster (Smaller) Builds
    \end{itemize}
\end{enumerate}

\subsection{Despliegue}
Subir carpeta Build/ a servidor con headers:
\begin{verbatim}
Content-Type: application/wasm (para .wasm)
Content-Encoding: br (si usa Brotli)
\end{verbatim}

\section{Extensión del Sistema}

\subsection{Añadir Nueva Pieza}
\begin{enumerate}
    \item Crear DronePartData asset (Create $>$ WebGL $>$ Drone Part Data)
    \item Configurar todos los campos
    \item Añadir ExplodablePart component al modelo 3D
    \item Asignar DronePartData al component
    \item Añadir ConnectionPointMarker a puntos de conexión
\end{enumerate}

\subsection{Añadir Nuevo Modo de Vista}
\begin{enumerate}
    \item Añadir entrada a enum ViewMode
    \item Crear material en ViewModeManager
    \item Implementar lógica en ApplyModeToRenderers()
    \item Añadir botón en ViewModeToolbar
\end{enumerate}

\subsection{Añadir Nuevo Manager}
\begin{enumerate}
    \item Heredar de Singleton<T> o PersistentSingleton<T>
    \item Implementar lógica en Awake() y Start()
    \item Suscribirse a eventos relevantes del EventBus
    \item Añadir al GameObject \_GameManagers en escena
\end{enumerate}

\section{Resumen de Scripts}

\begin{longtable}{|p{5cm}|p{8cm}|}
\hline
\textbf{Script} & \textbf{Función} \\
\hline
\endfirsthead
\hline
\textbf{Script} & \textbf{Función} \\
\hline
\endhead
AppStateMachine & Máquina de estados principal \\
EventBus & Sistema de eventos \\
SelectionManager & Selección de piezas \\
ViewModeManager & Modos de visualización (7) \\
PartCatalogManager & Búsqueda y filtrado \\
PartVisibilityManager & Ocultar/mostrar piezas \\
CrossSectionManager & Cortes transversales \\
DroneStateController & On/Off del dron \\
ModularPartsSystem & Intercambio de piezas \\
AssemblyGuideManager & Guía de ensamblaje \\
MeasurementTool & Mediciones 3D \\
ConnectionPointsViewer & Puntos de conexión \\
BillOfMaterialsManager & Lista de materiales \\
AnnotationSystem & Notas 3D \\
AssemblyChecklist & Verificación de piezas \\
AudioManager & Control de audio \\
TweenEngine & Motor de animaciones \\
OrbitCameraController & Cámara orbital \\
CameraPresets & Vistas predefinidas \\
HighlightSystem & Resaltado visual \\
CursorManager & Cursores dinámicos \\
LoadingController & Pantalla de carga \\
AnalyticsManager & Métricas de uso \\
AccessibilityManager & Opciones de accesibilidad \\
ErrorHandler & Manejo de errores \\
ScreenshotManager & Captura de pantalla \\
KeyboardShortcuts & Atajos de teclado \\
PerformanceMonitor & Métricas de rendimiento \\
WebGLOptimizer & Optimizaciones WebGL \\
\hline
\end{longtable}

\textbf{Total: 65+ scripts, 9,000+ líneas de código}

\end{document}
